\chapter*{Sommario}
\addcontentsline{toc}{chapter}{Sommario}

Il sequenziamento dell'RNA a singola cellula ha trasformato lo studio dell'eterogeneità cellulare, ma i protocolli a goccia ampiamente adottati catturano solo l'estremità $3'$ di ciascun trascritto \cite{conte2024opportunities}. Questo vincolo rende difficile il rilevamento dell'uso differenziale dei trascritti (DTU), poiché i read terminali brevi mancano delle giunzioni di splicing interne necessarie a distinguere le isoforme. Inoltre i metodi DTU basati sulle proporzioni---compresi quelli estesi ai dati a singola cellula, che si fondano sull'aggregazione in pseudo-bulk---sacrificano la risoluzione a livello di cellula per recuperare la profondità di lettura necessaria \cite{gilis2021satuRn, patrick2020sierra}.

Questa tesi indaga se COTAN---un framework statistico che modella direttamente i conteggi nulli tramite tabelle di contingenza di presenza e assenza \cite{galfre2021cotan}---possa essere esteso dalla co-espressione a livello genico al rilevamento DTU a livello di trascritto. Evitando sia l'imputazione artificiale sia l'aggregazione in pseudo-bulk, il metodo adattato identifica coppie di isoforme mutuamente esclusive tramite test di indipendenza e soglie di contrasto bidirezionale tra i cluster cellulari.

È stata costruita una pipeline Nextflow riproducibile per ottenere matrici a livello di trascritto da dati grezzi $3'$-biased usando alevin-fry con mappatura identità trascritto-gene \cite{he2022alevin}. La pipeline è stata applicata a una miscela di linee cellulari umane (29.000 cellule) e a tessuto cerebrale murino (4.000 cellule). Il clustering a livello di trascritto ha recuperato con successo sottostrutture a grana fine coerenti con il Global Differentiation Index, e i principali candidati DTU---tra cui Meg3, Malat1 e TPM1---sono risultati coerenti con la biologia delle isoforme documentata \cite{joglekar2023words}.

Questi risultati dimostrano che statistiche su tabelle di contingenza robuste alla sparsità possono rilevare l'esclusività reciproca tra isoforme in dati a goccia senza il costo proibitivo del sequenziamento full-length. Sebbene la sensibilità resti limitata dal bias posizionale---che rende invisibili gli eventi di splicing nelle regioni centrali del trascritto---e una validazione ortogonale sia ancora necessaria, questo lavoro stabilisce un percorso fondato per un'analisi isoform-aware su dataset atlas su larga scala.
